\chapter{Manual d'instal·lació}
\label{ch:annex-install}

Aquest annex documenta els passos per a desplegar una còpia local i pre-producció de Synapse Notes. El repositori és públic a \url{https://github.com/sergi-izquierdo/synapse-notes} i la versió etiquetada per a l'entrega del TFG és \texttt{v1.0-tfg-entrega}.

\section*{Requisits del sistema}

\begin{itemize}
    \item Node.js $\geq$ 24.0 (LTS recomanat). El \texttt{package.json} ho fa explícit amb el camp \texttt{engines}.
    \item NPM $\geq$ 10 (ve amb Node 24).
    \item Compte a Supabase amb un projecte creat (free tier suficient).
    \item Compte a Anthropic Console amb una clau API generada.
    \item Compte a Google AI Studio amb una clau API generada.
    \item Opcional per a deploy: compte a Vercel.
\end{itemize}

\section*{Passos d'instal·lació}

\begin{enumerate}
    \item \textbf{Clonar el repositori}:
    \begin{verbatim}
git clone https://github.com/sergi-izquierdo/synapse-notes.git
cd synapse-notes
git checkout v1.0-tfg-entrega
    \end{verbatim}

    \item \textbf{Instal·lar dependències}:
    \begin{verbatim}
npm install
    \end{verbatim}

    \item \textbf{Configurar variables d'entorn}. Copiar \texttt{.env.example} a \texttt{.env.local} i omplir les variables descrites a la taula~\ref{tab:env-vars}.

    \item \textbf{Aplicar les migracions SQL al Supabase del projecte}. Les migracions estan a \texttt{supabase/migrations/}. Cal aplicar-les en ordre cronològic (segons el prefix del nom). Si es té el CLI de Supabase instal·lat:
    \begin{verbatim}
npx supabase db push
    \end{verbatim}

    \item \textbf{Configurar els OAuth providers al Supabase}. Activar Google i GitHub a \texttt{Auth $\rightarrow$ Providers} del Dashboard de Supabase, configurant els callbacks segons la documentació oficial de cada proveïdor.

    \item \textbf{Iniciar el servidor de desenvolupament}:
    \begin{verbatim}
npm run dev
    \end{verbatim}

    La UI estarà disponible a \texttt{http://localhost:3000}.
\end{enumerate}

\begin{table}[H]
    \centerfloat
    \footnotesize
    \caption{Variables d'entorn del projecte (\texttt{.env.local}).}
    \label{tab:env-vars}
    \begin{tabular}{@{}l>{\raggedright\arraybackslash}p{6.5cm}l@{}}
        \toprule
        Variable & Propòsit & Obligatori \\
        \midrule
        \texttt{NEXT\_PUBLIC\_SUPABASE\_URL} & URL del projecte Supabase & sí \\
        \texttt{NEXT\_PUBLIC\_SUPABASE\_ANON\_KEY} & Clau anònima de Supabase (segura per al client) & sí \\
        \texttt{SUPABASE\_SERVICE\_ROLE\_KEY} & Clau de service-role (només servidor) & sí \\
        \texttt{ANTHROPIC\_API\_KEY} & Clau de l'API d'Anthropic & sí \\
        \texttt{GOOGLE\_GENERATIVE\_AI\_API\_KEY} & Clau de Google AI Studio & sí \\
        \texttt{AI\_ENABLED} & Kill-switch global (D4). Default \texttt{true}. & no \\
        \texttt{AI\_ALLOWLIST\_ONLY} & Activa allowlist mode (D4). Default \texttt{false}. & no \\
        \texttt{AI\_ALLOWLIST\_USER\_IDS} & UIDs autoritzats (CSV) si allowlist actiu. & no \\
        \texttt{AI\_JUDGE\_ENABLED} & Kill-switch del filtre D3. Default \texttt{true}. & no \\
        \texttt{AI\_JUDGE\_THRESHOLD} & Llindar de confiança per a bloqueig D3. Default 0{,}7. & no \\
        \bottomrule
    \end{tabular}
\end{table}

\chapter{Manual d'ús}
\label{ch:annex-usage}

Aquest annex resumeix el flux d'usuari principal de Synapse Notes des de l'autenticació fins a la consulta amb el xat RAG, il·lustrat amb captures de la versió desplegada a producció.

\section*{Primer accés}

\begin{enumerate}
    \item Obrir la URL del desplegament (per defecte \url{https://synapse-notes.vercel.app}).
    \item Triar un proveïdor OAuth (Google o GitHub) i autoritzar el redirect cap a Synapse Notes.
    \item En cas que el sistema estigui en mode allowlist (D4) i l'usuari no estigui inclòs, les funcionalitats d'IA mostraran un missatge informatiu. Les notes seguiran funcionant però sense RAG ni summarise.
\end{enumerate}

\begin{figure}[H]
    \centering
    \includegraphics[width=\linewidth]{screenshots/login.png}
    \caption{Pàgina d'entrada amb autenticació OAuth (Google i GitHub) via Supabase Auth.}
    \label{fig:ui-login}
\end{figure}

\section*{Crear i organitzar notes}

\begin{enumerate}
    \item Al dashboard, la zona \emph{ComposeZone} té un camp opcional de \emph{títol} i un textarea per al cos. El títol millora la cerca semàntica perquè s'inclou a la generació de l'embedding.
    \item Mentre s'escriu, els tres sigils invoquen autocomplete: \texttt{[[} obre cerca de notes per a backlinks, \texttt{@} per a mencions estil-tag, \texttt{\#} filtra tags existents.
    \item En desar, el sistema genera automàticament un embedding (Gemini) per a la nota i suggereix tags via Haiku 4.5 amb el component \emph{TagSuggestionRow}. L'usuari pot acceptar o rebutjar.
    \item Star, archive, delete i restore són accions disponibles a cada card amb actualitzacions optimistes.
\end{enumerate}

\section*{Xat RAG}

\begin{enumerate}
    \item Obrir la barra lateral del xat (icona a la cantonada inferior-dreta a mòbil, integrada al panel esquerre a desktop).
    \item Escriure una consulta en llenguatge natural (qualsevol idioma; el sistema detecta i respon en el mateix).
    \item El sistema mostra els raonaments del model i invoca tools (\texttt{getNotesByTag}, \texttt{graph\_neighbors}, \texttt{graph\_shortest\_path}) quan és necessari.
    \item Les referències a notes apareixen com \texttt{[id=N]} o clicables amb el seu títol.
\end{enumerate}

\begin{figure}[H]
    \centering
    \includegraphics[width=0.5\linewidth]{screenshots/chat-rag.png}
    \caption{Xat RAG responent una consulta per etiqueta. El model invoca l'eina \texttt{getNotesByTag} i retorna la llista de notes amb el tag \texttt{Compra} i un resum del contingut de cadascuna.}
    \label{fig:ui-chat}
\end{figure}

\section*{Visualització del graph}

\begin{enumerate}
    \item Navegar a \texttt{/graph}.
    \item Al canvas 2D, els nodes són notes, les arestes representen relacions per tag compartit, similitud semàntica o backlink directe.
    \item Hover sobre un node mostra previsualització; click obre l'inspector lateral.
    \item Toggle 2D/3D a la part superior del side panel canvia entre les dues vistes. La 3D té rotació orbital amb el ratolí.
    \item La cerca textual filtra el graph en temps real, dimming els nodes no coincidents.
\end{enumerate}

\chapter{Connexió al servidor MCP via OAuth 2.1}
\label{ch:annex-mcp}

Aquest annex descriu com connectar un client MCP (Claude Code, Claude Desktop, o qualsevol agent compatible amb la especificació MCP) al servidor de Synapse Notes. El flux és OAuth~2.1 amb PKCE i Dynamic Client Registration (RFC~7591), sense necessitat d'extreure manualment cap token.

\section*{Pas 1: Registrar el servidor MCP al client}

A la terminal, executar:

\begin{verbatim}
claude mcp add synapse-notes \
  --transport http \
  https://synapse-notes.vercel.app/api/mcp
\end{verbatim}

Per a Claude Desktop, afegir l'entrada equivalent a \texttt{claude\_desktop\_config.json} amb \texttt{transport: "http"} i la mateixa URL. El client detecta automàticament que el servidor requereix OAuth (rep un \texttt{401} amb \texttt{WWW-Authenticate} que apunta a \texttt{/.well-known/oauth-protected-resource}).

\section*{Pas 2: Autoritzar via navegador}

El client obre automàticament el navegador a la pàgina de consentiment \texttt{/mcp-authorize}. La pàgina:

\begin{enumerate}
    \item Verifica que l'usuari té sessió activa a Synapse Notes (si no, redirigeix al login OAuth Google/GitHub).
    \item Mostra el nom del client sol·licitant i els permisos demanats (read/write de notes, gestió de tags, sumarització).
    \item Presenta els botons \emph{Authorize} / \emph{Deny}.
\end{enumerate}

Al clicar \emph{Authorize}, el servidor genera un codi d'autorització xifrat (AES-256-GCM, TTL de 5~minuts) que encapsula els tokens Supabase de la sessió. El client el rep via redirect i l'intercanvia per un \texttt{access\_token} + \texttt{refresh\_token} contra l'endpoint \texttt{/api/oauth/token}, verificant PKCE~S256.

\section*{Pas 3: Verificar la connexió}

A Claude Code, executar \texttt{/mcp} i comprovar que \texttt{synapse-notes} apareix com a \emph{authenticated}. Provar una crida:

\begin{quote}
``Cerca a les meves notes la paraula \emph{tutorial}.''
\end{quote}

El client invocarà \texttt{search\_notes} via MCP, retornarà els resultats i els resumirà. La renovació del token quan caduca és transparent: el client crida \texttt{/api/oauth/token} amb \texttt{grant\_type=refresh\_token} i Supabase emet un nou \texttt{access\_token} sense intervenció de l'usuari.

\paragraph{Diferència amb el flux manual original.} La versió inicial del servidor MCP requeria extreure manualment el JWT de Supabase del cookie del navegador i enganxar-lo a la configuració. Aquesta solució tenia tres problemes: el token quedava en clar al disc del client, caducava en 1~hora sense renovació automàtica, i donava accés complet a la sessió de l'usuari sense scope. El flux OAuth~2.1 d'aquest annex resol els tres problemes seguint l'estàndard MCP-spec (vegeu \S\ref{sec:impl-oauth-mcp} per a la implementació detallada).

\chapter{Repositori i artefactes del projecte}
\label{ch:annex-repo}

\section*{Codi font}

\begin{itemize}
    \item \textbf{Repositori principal}: \url{https://github.com/sergi-izquierdo/synapse-notes}
    \item \textbf{Tag d'entrega}: \texttt{v1.0-tfg-entrega} (es crearà al moment de l'entrega).
    \item \textbf{Branca de desenvolupament}: \texttt{main}.
\end{itemize}

\section*{Artefactes verificables}

\begin{itemize}
    \item \texttt{docs/tfg/00-decision-log.md}: 5 decisions arquitectòniques majors (D1--D5) amb segells temporals.
    \item \texttt{docs/tfg/extend.md}: log granular de progrés setmana a setmana.
    \item \texttt{promptfoo/}: suite red-team amb 36 variants i scripts d'execució + cleanup. Fitxers de resultats raw a \texttt{promptfoo/results/}.
    \item \texttt{supabase/migrations/}: tota la història d'evolució de l'schema des de la versió MCP.
    \item \texttt{graphify-out/graph.html}: graph interactiu del propi repositori (auditoria estructural via graphify del 2026-04-24).
\end{itemize}

\section*{Reproducció del build del memoir}

El memoir està escrit en LaTeX (memoir class) amb \texttt{lualatex} + \texttt{biber}. Per a reproduir el PDF:

\begin{verbatim}
cd tfg/
mmdc -i figures/c4-context.mmd -o figures/c4-context.pdf
mmdc -i figures/c4-containers.mmd -o figures/c4-containers.pdf
mmdc -i figures/mcp-oauth-sequence.mmd -o figures/mcp-oauth-sequence.pdf

lualatex main.tex
biber main
lualatex main.tex
lualatex main.tex
\end{verbatim}

El procés requereix MiKTeX o TeX Live amb els paquets memoir, biblatex-apa, pgfplots, fontspec i graphicx. La regeneració de les figures Mermaid requereix Node.js i \texttt{@mermaid-js/mermaid-cli} (\texttt{mmdc}).

\section*{Demo i materials suplementaris}

\begin{itemize}
    \item \textbf{URL de la demo}: \url{https://synapse-notes.vercel.app} (sotmesa a allowlist D4).
    \item \textbf{Vídeo de demostració}: enllaç pendent (es generarà al moment de l'entrega).
    \item \textbf{Slides de defensa}: pendents (es prepararan setmana 7 segons el plan).
\end{itemize}
